\documentclass[11pt]{amsart}
\usepackage{amsmath,amssymb,amsthm}
\usepackage[margin=1.1in]{geometry}
\usepackage{hyperref}

\newtheorem{theorem}{Theorem}
\newtheorem{corollary}[theorem]{Corollary}
\newtheorem{lemma}[theorem]{Lemma}
\theoremstyle{remark}
\newtheorem{remark}[theorem]{Remark}

\newcommand{\F}{\mathbb{F}_2}
\newcommand{\disc}{\operatorname{disc}}
\newcommand{\NL}{\operatorname{NL}}

\title{The discrepancy of Sylvester--Hadamard matrices in odd dimension}
\author{Yitzchak Shmalo}
\date{}

\begin{document}

\begin{abstract}
Let $H_k$ be the $2^k\times 2^k$ Sylvester--Hadamard matrix and let
$\disc(H_k)$ be the minimum of $\|H_kx\|_\infty$ over sign vectors $x$. It is known that
$\disc(H_k)=2^{k/2}$ for even $k$, and it has been conjectured that
$\disc(H_k)=2^{(k+1)/2}$ for every odd $k$. We show that
$\disc(H_k)=2^k-2\rho(1,k)$, where $\rho(1,k)$ is the covering radius of the
first-order Reed--Muller code of length $2^k$. Known values and bounds for this
covering radius then show that the conjecture holds for $k\le 7$ and fails for
every odd $k\ge 9$; in fact $\disc(H_k)\le \tfrac78\,2^{(k+1)/2}$ for every odd
$k\ge 9$. Determining $\disc(H_k)$ for odd $k\ge 9$ is therefore equivalent to
determining $\rho(1,k)$.
\end{abstract}

\maketitle

\section{Introduction}

For a real $m\times n$ matrix $A$ the discrepancy is
\[
\disc(A)=\min_{x\in\{-1,1\}^n}\|Ax\|_\infty .
\]
Index the rows and columns of the Sylvester--Hadamard matrix $H_k$ by $\F^k$,
so that $(H_k)_{a,y}=(-1)^{a\cdot y}$. By Parseval every $x\in\{-1,1\}^{2^k}$
satisfies $\|H_kx\|_2^2=2^{2k}$, hence $\|H_kx\|_\infty\ge 2^{k/2}$, and
$\disc(H_k)=2^{k/2}$ when $k$ is even. The problem list \cite{RS101} records the
odd case as open and states the following.

\medskip
\noindent\textbf{Conjecture} (\cite[Conjecture 13]{RS101}).
For every odd $k$, $\disc(H_k)=\sqrt2\,\sqrt{2^k}=2^{(k+1)/2}$.
\medskip

The value $2^{(k+1)/2}$ is what one obtains from $H_{k-1}$ and a bent function,
and the question is whether it can be improved. The answer is a statement about
Boolean functions that is classical in coding theory. For $f\colon\F^k\to\F$ let
\[
W_f(a)=\sum_{y\in\F^k}(-1)^{f(y)+a\cdot y}
\]
be its Walsh transform, and let $\NL(f)$ be the Hamming distance from $f$ to the
set of affine functions. The covering radius of the first-order Reed--Muller code
$RM(1,k)$, the code of affine functions, is $\rho(1,k)=\max_f\NL(f)$.

\begin{theorem}\label{thm:main}
For every $k\ge1$, $\disc(H_k)=2^k-2\rho(1,k)$.
\end{theorem}

\begin{corollary}\label{cor}
\begin{enumerate}
\item $\disc(H_k)=2^{(k+1)/2}$ for $k\in\{1,3,5,7\}$.
\item $\disc(H_9)\le 28<32=2^{5}$.
\item For every odd $k\ge9$, $\disc(H_k)\le \tfrac78\,2^{(k+1)/2}$. In particular the
Conjecture fails for every odd $k\ge 9$.
\item For every odd $k\ge15$, $\disc(H_k)\le \tfrac{27}{32}\,2^{(k+1)/2}$.
\end{enumerate}
\end{corollary}

\section{Proofs}

\begin{proof}[Proof of Theorem~\ref{thm:main}]
Every $x\in\{-1,1\}^{\F^k}$ is $x=(-1)^f$ for a unique $f\colon\F^k\to\F$, and then
$(H_kx)_a=\sum_y(-1)^{a\cdot y+f(y)}=W_f(a)$. Writing $d(f,g)$ for Hamming
distance,
\[
d(f,a\cdot y)=\tfrac12\bigl(2^k-W_f(a)\bigr),\qquad
d(f,a\cdot y+1)=\tfrac12\bigl(2^k+W_f(a)\bigr),
\]
so the distance from $f$ to the affine functions is
$\NL(f)=\tfrac12\bigl(2^k-\max_a|W_f(a)|\bigr)$. Hence
\[
\|H_kx\|_\infty=\max_a|W_f(a)|=2^k-2\NL(f),
\]
and minimizing over $x$ is maximizing $\NL(f)$ over $f$.
\end{proof}

\begin{lemma}\label{lem:sum}
Let $f\colon\F^k\to\F$ and let $b\colon\F^m\to\F$ be bent ($m$ even). The function
$g(y,z)=f(y)+b(z)$ on $\F^{k+m}$ satisfies $\max|W_g|=2^{m/2}\max|W_f|$.
Consequently $\disc(H_{k+m})\le 2^{m/2}\disc(H_k)$ for every even $m\ge0$.
\end{lemma}

\begin{proof}
The sum defining $W_g(a,c)$ factors as $W_f(a)\,W_b(c)$, and $|W_b(c)|=2^{m/2}$
for every $c$ because $b$ is bent. Bent functions exist in every even
dimension \cite{Rothaus}. The second statement follows from
Theorem~\ref{thm:main}, applied with $f$ of maximal nonlinearity.
\end{proof}

\begin{proof}[Proof of Corollary~\ref{cor}]
(1) The covering radii are $\rho(1,1)=0$ and $\rho(1,3)=2$ by inspection,
$\rho(1,5)=12$ \cite{BW}, and $\rho(1,7)=56$ \cite{Mykkeltveit}. Theorem~\ref{thm:main}
gives $\disc(H_k)=2,4,8,16$ for $k=1,3,5,7$, which is $2^{(k+1)/2}$ in each case.

(2) Kavut and Y\"ucel \cite{KY} found nine-variable Boolean functions of
nonlinearity $242$, so $\rho(1,9)\ge242$ and
$\disc(H_9)\le 512-484=28$. (Nonlinearity $241$, found earlier by Kavut, Maitra
and Y\"ucel \cite{KMY}, already gives $\disc(H_9)\le30$.)

(3) Every odd $k\ge9$ is $9+m$ with $m$ even. By Lemma~\ref{lem:sum} and (2),
$\disc(H_k)\le 28\cdot 2^{m/2}=\tfrac{28}{32}\,2^{(9+m+1)/2}$.

(4) Patterson and Wiedemann \cite{PW} showed $\rho(1,15)\ge16276$, so
$\disc(H_{15})\le 32768-32552=216=\tfrac{27}{32}\cdot 256$. Lemma~\ref{lem:sum}
carries this to every odd $k\ge15$.
\end{proof}

\section{Remarks}

\begin{remark}
Part (4) uses only results available since 1983, so the Conjecture was false for
$k\ge15$ before it was stated. The first odd dimension in which it fails is $k=9$.
\end{remark}

\begin{remark}
By Theorem~\ref{thm:main} the odd case of the problem in \cite{RS101} is the
problem of the covering radius of $RM(1,k)$ for odd $k$, open since the 1970s.
For odd $k$ the known bounds are
\[
2^{k/2}<\disc(H_k)\le 2^{(k+1)/2},
\]
where the strict lower bound is the classical fact that bent functions exist only in
even dimension, and the upper bound is improved by the constructions above for every
odd $k\ge9$. Any new lower bound on $\rho(1,k)$, and any upper bound on
$\rho(1,k)$, transfers to $\disc(H_k)$ through Theorem~\ref{thm:main} without loss.
\end{remark}

\begin{remark}
The list \cite{RS101} also states Conjecture 12,
$\limsup_{n\to\infty}\sup_{A\in\{\pm1\}^{n\times n}}\disc(A)/\sqrt n>1$, and
notes that Conjecture 13 would imply it through $H_k$, $k$ odd. By
Theorem~\ref{thm:main}, Sylvester--Hadamard matrices witness Conjecture 12 if and
only if
\[
\limsup_{k\ \mathrm{odd}}\ \bigl(2^k-2\rho(1,k)\bigr)\,2^{-k/2}>1 .
\]
Corollary~\ref{cor} bounds this $\limsup$ by $\tfrac{27}{32}\sqrt2\approx1.19$; whether it
exceeds $1$ is a question about the covering radius of $RM(1,k)$ that the present
note leaves open.
\end{remark}

\begin{remark}
The repository that accompanies this note contains \texttt{check.py}, which
verifies Theorem~\ref{thm:main} by exhaustive search for $k\le4$ and checks the
product identity of Lemma~\ref{lem:sum} on random functions.
\end{remark}

\begin{thebibliography}{9}

\bibitem{RS101}
Randomstrasse101, \emph{Open problems of 2024}, arXiv:2504.20539, 2025.

\bibitem{BW}
E.~R.~Berlekamp and L.~R.~Welch, Weight distributions of the cosets of the
$(32,6)$ Reed--Muller code, \emph{IEEE Trans. Inform. Theory} 18 (1972), 203--207.

\bibitem{Mykkeltveit}
J.~J.~Mykkeltveit, The covering radius of the $(128,8)$ Reed--Muller code is $56$,
\emph{IEEE Trans. Inform. Theory} 26 (1980), 359--362.

\bibitem{PW}
N.~J.~Patterson and D.~H.~Wiedemann, The covering radius of the $(2^{15},16)$
Reed--Muller code is at least $16276$, \emph{IEEE Trans. Inform. Theory} 29 (1983),
354--356.

\bibitem{KMY}
S.~Kavut, S.~Maitra and M.~D.~Y\"ucel, Search for Boolean functions with excellent
profiles in the rotation symmetric class, \emph{IEEE Trans. Inform. Theory} 53 (2007),
1743--1751.

\bibitem{KY}
S.~Kavut and M.~D.~Y\"ucel, $9$-variable Boolean functions with nonlinearity $242$
in the generalized rotation symmetric class, \emph{Inform. and Comput.} 208 (2010),
341--350.

\bibitem{Rothaus}
O.~S.~Rothaus, On ``bent'' functions, \emph{J. Combin. Theory Ser. A} 20 (1976),
300--305.

\end{thebibliography}

\end{document}
